\documentclass[11pt,a4paper]{article}
\usepackage[margin=1in]{geometry}
\usepackage[hidelinks]{hyperref}
\usepackage{enumitem}
\usepackage{graphicx}
\usepackage{xcolor}
\usepackage{amsmath}

% -----------------------------------------------------
% bvraghav-tiet-ucs503p-article.sty
% -----------------------------------------------------

% Variable: Lab Instructor
\makeatletter \newcommand{\BvrSetLabInstructor}[1]{%
  \gdef\Bvr@LabInstructor{#1}}
% default
\newcommand{\Bvr@LabInstructor}{Dr. Raghav %
  B. Venkataramaiyer}
% public accessor
\newcommand{\BvrLabInstructorValue}{\Bvr@LabInstructor}
\makeatother

% Add subtitle
\usepackage{titling}
% -- Set title bold
\pretitle{\begin{center}\bfseries\fontsize{18bp}{18bp}\selectfont}
\makeatletter
\newcommand{\BvrSubtitle}[1]{%
  \posttitle{%
    \par\end{center}
    \begin{center}\large#1\end{center}
    \vskip 0.5em}%
}
\makeatother

% Hints in the section title(s)
\newcommand{\BvrHint}[1]{%
  {\normalfont\normalsize\itshape\hfill #1}}

% -----------------------------------------------------

\title{Cardio AI: Explainable Heart Disease Prediction using\\CNN-LSTM Attention Network with Grad-CAM Interpretability}

\BvrSetLabInstructor{Nisha Thakur}

\BvrSubtitle{%
  Project Proposal for UCS503P  \\
  Submitted to: \BvrLabInstructorValue \\ \bfseries
  Thapar Institute of Engineering and Technology%
}

\author{
  Dhruv Singla, CSED%
  \thanks{Roll No: \texttt{1024170434}, %
    Email: \texttt{dsingla3\_be24@thapar.edu}}
  \and
  Srishti, CSED%
  \thanks{Roll No: \texttt{1024170451}, %
    Email: \texttt{ssrishti1\_be24@thapar.edu}}
}

\date{\today}

\begin{document}
\maketitle

% ===========================================================================
\section{Higher-order goal%
  \BvrHint{\ldots secondary framing}}
% ===========================================================================
This project aims to improve \emph{clinical transparency in cardiac
diagnostics} by making AI-driven heart disease predictions
interpretable to both clinicians and patients. While the broader
aspiration is trustworthy medical AI, the immediate engineering
objective is to build a deployable, explainable prediction system
that translates deep learning confidence scores into actionable
clinical narratives.

% ===========================================================================
\section{Time-to-value%
  \BvrHint{\ldots primary heuristic}}
% ===========================================================================
\begin{itemize}[leftmargin=0.7cm]
    \item We will deliver meaningful increments quickly by first
      training and validating the CNN-LSTM-Attention model on the
      Cleveland dataset, then layering explainability (Grad-CAM,
      Attention, SHAP) and the clinical web interface in successive
      iterations.
    \item We will prioritize fast feedback over perfection by using
      short iteration cycles: model accuracy is validated before
      explainability modules are integrated, and each module is
      independently testable.
    \item Success is defined by what can be measured in the first
      iteration (holdout AUC $\geq 0.85$), then improved based on
      cross-validation and clinical-alignment feedback.
\end{itemize}

% ===========================================================================
\section{Problem Statement}
% ===========================================================================
Cardiovascular disease remains the leading cause of mortality
globally. Early prediction from routine clinical measurements
(blood pressure, cholesterol, ECG markers) can significantly
improve patient outcomes. However, current challenges include:
\begin{itemize}[leftmargin=0.7cm]
    \item \textbf{Black-box predictions:} Deep learning models
      achieve high accuracy but offer no clinical reasoning, making
      physicians reluctant to trust automated diagnoses.
    \item \textbf{Feature opacity:} With 13 heterogeneous clinical
      features (numerical vitals, categorical ECG types, binary
      indicators), understanding \emph{which} features drive a
      specific patient's risk is non-trivial.
    \item \textbf{Limited visual evidence:} Clinicians lack
      visual artifacts (heatmaps, feature importance charts) that
      map model confidence back to physiological parameters.
    \item \textbf{No longitudinal tracking:} Existing tools
      rarely support patient-history-aware predictions or trend
      analysis across repeated assessments.
\end{itemize}

\textbf{Why this matters now (contextual relevance):} As AI-assisted
diagnosis tools enter clinical workflows, regulatory bodies and
medical ethics demand interpretable, auditable predictions ---
not just accurate ones.

% ===========================================================================
\section{Proposed Solution%
  \BvrHint{\ldots Software Proposition}}
% ===========================================================================

\subsection{Overview}
Build a \textbf{web-based} explainable cardiac risk prediction
system with the following components:
\begin{itemize}[leftmargin=0.7cm]
    \item \textbf{CNN-LSTM-Attention deep learning model:} trained
      on the UCI Cleveland Heart Disease dataset (303 records,
      13 clinical features) with GAF image encoding for
      spatiotemporal feature extraction.
    \item \textbf{Multi-method explainability engine:} Grad-CAM
      heatmaps, attention weight visualization, and SHAP values
      providing three complementary lenses on model reasoning.
    \item \textbf{Clinical report generator:} automated PDF reports
      with narrative summaries, feature-by-feature clinical
      interpretation, and guideline-referenced evaluations.
    \item \textbf{Patient management portal:} Flask-based web
      interface with patient registration, longitudinal history
      tracking, and searchable records via MySQL backend.
    \item \textbf{Interactive visualization dashboard:} Plotly-based
      charts for real-time XAI comparison (Attention vs.\ Grad-CAM
      vs.\ SHAP).
\end{itemize}

\subsection{Core workflow}
\begin{enumerate}[leftmargin=0.7cm]
    \item Clinician enters patient clinical parameters (age, blood
      pressure, cholesterol, ECG markers, etc.) via the web form.
    \item System preprocesses inputs using a ClinicalRangeScaler
      (physiological-bound normalization), converts to GAF image
      representation, and runs prediction through the
      CNN-LSTM-Attention model.
    \item Explainability engine generates Grad-CAM heatmap,
      extracts attention weights, and computes SHAP attributions.
    \item System produces an interactive dashboard with XAI
      visualizations and a downloadable PDF clinical report with
      guideline-referenced narrative.
    \item Patient record is persisted to MySQL database for
      longitudinal tracking and future comparisons.
\end{enumerate}

\subsection{Operational constraints for an educational setting}
\begin{itemize}[leftmargin=0.7cm]
    \item Use the publicly available Cleveland dataset (303 records)
      --- no patient consent or privacy concerns.
    \item Keep the web interface responsive and usable on standard
      hardware without GPU requirements at inference time.
    \item Design for local deployment (Flask + MySQL) with a clear
      path to containerized hosting.
\end{itemize}

% ===========================================================================
\section{Solution Approach%
  \BvrHint{\ldots Engineering focus}}
% ===========================================================================
This project combines deep learning engineering with clinical
interpretability, focusing on robust pipeline design and
multi-method XAI rather than novel architecture research.

\subsection{Data preprocessing pipeline%
  \BvrHint{\ldots non-research}}
\begin{itemize}[leftmargin=0.7cm]
    \item \textbf{Missing value handling:} median imputation for
      Cleveland \texttt{ca} and \texttt{thal} columns (marked as
      \texttt{?}).
    \item \textbf{Outlier treatment:} IQR-based capping to preserve
      dataset size on a small ($n=303$) dataset.
    \item \textbf{Feature scaling:} custom ClinicalRangeScaler using
      fixed physiological boundaries (e.g., age: 25--85, BP:
      80--200~mmHg) instead of dataset-dependent statistics,
      ensuring out-of-sample consistency.
    \item \textbf{Encoding:} one-hot encoding for categorical
      features (chest pain type, thal, slope, etc.) with explicit
      feature-to-index mapping for XAI aggregation.
    \item \textbf{Image transformation:} Gramian Angular Summation
      Field (GASF) encoding converts 1D feature vectors to 2D
      images for Conv2D input.
    \item \textbf{Class balancing:} SMOTE oversampling on the
      training split with fallback to class-weighted loss.
\end{itemize}

\subsection{Model architecture%
  \BvrHint{\ldots deep learning}}
A hybrid CNN-LSTM-Attention architecture:
\begin{itemize}[leftmargin=0.7cm]
    \item \textbf{Conv2D layers:} extract spatial correlations from
      GAF-encoded feature images.
    \item \textbf{LSTM layer:} captures sequential/temporal
      dependencies across feature dimensions.
    \item \textbf{Attention mechanism:} learns to weight feature
      contributions, providing an inherently interpretable
      bottleneck.
    \item \textbf{Output:} sigmoid activation for binary
      classification (heart disease present/absent).
\end{itemize}

\subsection{Explainability methods%
  \BvrHint{\ldots interpretable AI}}
Three complementary XAI techniques:
\begin{itemize}[leftmargin=0.7cm]
    \item \textbf{Grad-CAM:} gradient-weighted class activation maps
      on the final Conv2D layer, averaged to 1D feature importance.
    \item \textbf{Attention weights:} extracted from the attention
      layer and aggregated from one-hot indices back to original
      clinical features.
    \item \textbf{SHAP values:} model-agnostic Shapley attributions
      providing per-feature contribution with directionality.
\end{itemize}

\subsection{Web architecture%
  \BvrHint{\ldots web-based solution}}
A 3-tier architecture:
\begin{itemize}[leftmargin=0.7cm]
    \item \textbf{Frontend:} HTML/CSS/JS with Plotly.js for
      interactive XAI charts and responsive patient input forms.
    \item \textbf{Backend API:} Flask (Python) serving prediction,
      explainability, PDF generation, and patient search endpoints.
    \item \textbf{Database:} MySQL storing patients, clinical
      records, AI diagnoses, and PDF report paths.
\end{itemize}

% ===========================================================================
\section{Evaluation Criterion%
  \BvrHint{\ldots measurable and attributable}}
% ===========================================================================
We define success using measurable metrics tied to both model
performance and clinical interpretability.

\subsection{Primary evaluation metric}
\textbf{Area Under the ROC Curve (AUC-ROC):} measures the
model's ability to discriminate between heart disease and
normal cases across all classification thresholds.
\begin{itemize}[leftmargin=0.7cm]
    \item Target: AUC $\geq 0.85$ on the 20\% holdout test set.
    \item Attribution: computed from model probability outputs
      against ground-truth binary labels.
\end{itemize}

\subsection{Secondary evaluation metrics}
\begin{itemize}[leftmargin=0.7cm]
    \item \textbf{Accuracy, Precision, Recall, F1-Score:}
      standard classification metrics on the holdout test set.
    \item \textbf{Cross-validation stability:} 5-fold stratified
      CV to verify generalization (target: CV AUC std $\leq 0.05$).
    \item \textbf{XAI clinical alignment:} top-3 features
      identified by Grad-CAM, Attention, and SHAP are
      cross-referenced against established clinical risk factors
      (e.g., \texttt{ca}, \texttt{thal}, \texttt{cp},
      \texttt{oldpeak}) per ACC/AHA guidelines.
    \item \textbf{Report completeness:} generated PDF reports
      contain narrative summary, feature-level evaluation against
      clinical guidelines, and risk-stratified recommendations.
\end{itemize}

\subsection{Pilot validation plan%
  \BvrHint{\ldots high-level}}
\begin{itemize}[leftmargin=0.7cm]
    \item Validate predictions on all 303 Cleveland dataset records
      using nearest-neighbour ground-truth matching.
    \item Demonstrate the web portal end-to-end: form submission
      $\rightarrow$ prediction $\rightarrow$ XAI dashboard
      $\rightarrow$ PDF report $\rightarrow$ patient history
      retrieval.
    \item Collect qualitative feedback from 3--5 peers on report
      readability and XAI visualization clarity.
\end{itemize}

% ===========================================================================
\section{Scalability%
  \BvrHint{\ldots with foresight}}
% ===========================================================================
The proposition should scale in both theoretical and practical senses.

\begin{enumerate}[leftmargin=0.7cm]
    \item \textbf{Scalable (theoretically):} The system should
    support growth in patient volume and model complexity by:
    \begin{itemize}[leftmargin=0.7cm]
        \item decoupling the prediction engine from the web server
          (API-first design),
        \item stateless API design supporting future horizontal
          scaling,
        \item using SQL queries for patient history search.
    \end{itemize}

    \item \textbf{Operationally deployable (practically):} The
    system should run on standard infrastructure:
    \begin{itemize}[leftmargin=0.7cm]
        \item Flask + Docker-hosted MySQL on any Windows/Linux/macOS host,
        \item containerizable via Docker for reproducible deployment,
        \item model weights serialized for inference without
          retraining.
    \end{itemize}
\end{enumerate}

% ===========================================================================
\section{Engine availability heuristic}
% ===========================================================================
A mature set of ``engines'' (tooling/libraries) is available to
implement this project:
\begin{itemize}[leftmargin=0.7cm]
    \item \textbf{TensorFlow/Keras:} CNN-LSTM-Attention model
      construction, training, Grad-CAM computation.
    \item \textbf{scikit-learn:} preprocessing, train-test splitting,
      classification metrics, cross-validation.
    \item \textbf{imbalanced-learn:} SMOTE for class balancing.
    \item \textbf{SHAP:} model-agnostic Shapley value computation.
    \item \textbf{Flask:} lightweight web framework for REST API
      and template rendering.
    \item \textbf{PyMySQL:} MySQL database connectivity.
    \item \textbf{Plotly.js:} interactive client-side XAI
      visualizations.
    \item \textbf{Microsoft Edge (headless):} HTML-to-PDF clinical
      report conversion.
    \item \textbf{Pandas, NumPy, Matplotlib, Seaborn:} data
      manipulation, EDA, and static plotting.
\end{itemize}
We will use these mature components to focus on pipeline
correctness, explainability quality, and clinical alignment rather
than inventing new infrastructure.

% ===========================================================================
\section{Project scope and deliverables%
  \BvrHint{\ldots iteration-friendly}}
% ===========================================================================

\subsection{Initial deliverable%
  \BvrHint{\ldots first iteration}}
\begin{itemize}[leftmargin=0.7cm]
    \item EDA pipeline: missing values, duplicates, class
      distribution, correlation analysis
    \item Preprocessing pipeline: ClinicalRangeScaler, one-hot
      encoding, GASF transformation, SMOTE
    \item CNN-LSTM-Attention model: training with holdout evaluation
      and 5-fold cross-validation
    \item Grad-CAM and Attention weight extraction with
      one-hot-to-original feature aggregation
    \item Basic Flask endpoint: accepts patient features, returns
      prediction and XAI weights
\end{itemize}

\subsection{Subsequent deliverables}
\begin{itemize}[leftmargin=0.7cm]
    \item SHAP integration for model-agnostic feature attribution
    \item Interactive Plotly.js dashboard for XAI comparison
      (Attention vs.\ Grad-CAM vs.\ SHAP)
    \item PDF clinical report generator with guideline-referenced
      narrative summaries
    \item Patient management: registration, longitudinal history
      tracking, search API
    \item GAF image visualization and 2D Grad-CAM heatmap display
    \item Subgroup-aware performance metrics (by age group, sex)
\end{itemize}

% ===========================================================================
\section{Risks and mitigations}
% ===========================================================================
\begin{itemize}[leftmargin=0.7cm]
    \item \textbf{Small dataset ($n=303$):} mitigate via
      stratified splitting, SMOTE, cross-validation, and IQR
      outlier capping to maximize data utility.
    \item \textbf{Overfitting on small data:} EarlyStopping,
      ReduceLROnPlateau, and ModelCheckpoint callbacks; dropout
      regularization in the network.
    \item \textbf{XAI disagreement across methods:} present all
      three methods (Grad-CAM, Attention, SHAP) side-by-side;
      highlight consensus features rather than relying on a single
      method.
    \item \textbf{Clinical misinterpretation:} clearly label all
      outputs as ``AI-assisted, not diagnostic''; reference
      ACC/AHA clinical guidelines in reports.
    \item \textbf{TensorFlow deployment issues:} maintain a
      lightweight fallback inference path using NumPy-based
      forward pass for environments without GPU/native TF support.
\end{itemize}

% ===========================================================================
\section{Summary}
% ===========================================================================
This project proposes a deployable, explainable heart disease
prediction system combining a CNN-LSTM-Attention deep learning
model with multi-method interpretability (Grad-CAM, Attention
weights, SHAP). It meets the course criteria by emphasizing
time-to-value through iterative model--explainability--interface
delivery, using measurable evaluation criteria (AUC-ROC $\geq 0.85$,
cross-validation stability, clinical feature alignment), and
producing clinician-readable PDF reports. The system is designed
for scalable deployment via a Flask web portal with MySQL-backed
patient history management.

\end{document}
